\subsection{Acoustic Dataset: FSC22}

The acoustic threat classifier is trained on a curated subset of the FSC22
(Forest Sound Classification 2022) dataset \cite{kahl2023fsc22}. FSC22
contains 2,025 five-second audio clips across 27 sub-classes grouped under
six broader categories, each holding 75 samples drawn from the Freesound
library.

For this project, FSC22 was narrowed to six project-relevant classes: Axe,
Chainsaw, Gunshot, Handsaw, Tree Falling, and a merged Background class
combining twelve non-threat sub-classes (rain, wind, stream, birds, insects,
footsteps, and others). Table~\ref{tab:class-mapping} shows this mapping.
All non-forest classes (Car, Train, Engine, Motorcycle, Helicopter, etc.)
were excluded as out of scope.

\begin{table}[H]
    \centering
    \caption{FSC22 sub-class to project-class mapping}
    \label{tab:class-mapping}
    \begin{tabular}{ll}
        \toprule
        \textbf{Project Class} & \textbf{Mapped FSC22 Sub-classes} \\
        \midrule
        Axe & Axe, Wood\_Chop \\
        Chainsaw & Chainsaw, Sawing, Wood\_Sawing \\
        Gunshot & Gunshot, Gunshot\_Other \\
        Handsaw & Handsaw, Sawing, Wood\_Sawing \\
        Tree Falling & Tree\_Falling, Wood\_Crack, Wood\_Creak \\
        Background & Rain, Wind, Stream, Ocean, Crackling\_Fire, Steps, \\
                    & Footwalk, Bird, Bird\_Other, Insects, Thunder, Church\_Bell \\
        \bottomrule
    \end{tabular}
\end{table}

This 6-class formulation was used for the Full and Compact architecture
designs described below (Section~\ref{subsec:model-evolution}). During
training and evaluation of the deployed Tiny architecture, the Tree Falling
class was dropped: its low-frequency crack/creak signature overlapped
heavily with background forest noise and its removal improved the
reliability of the remaining classes. The final Tiny model is therefore
trained and evaluated on five classes -- Axe, Chainsaw, Gunshot, Handsaw,
and Background -- and all classification results reported in
Section~\ref{subsec:training-pipeline} and Section~\ref{sec:acoustic-results}
reflect this 5-class configuration.

\subsection{Acoustic Feature Extraction}

Raw 16~kHz audio is converted to a log-Mel spectrogram via a Hann window,
FFT, and Mel filterbank, followed by per-frame normalisation. This
time-frequency representation captures the harmonic content of chainsaw
noise, the sharp transients of gunshots, and the impact profile of axe
strikes.

\subsection{Memory-Constrained Model Architecture Design}
\label{subsec:model-evolution}

Designing a CNN for on-device inference on the ESP32-S3 required an
iterative architecture search, since the peak intermediate tensor size at
each convolution layer must fit within the chip's 512~KB SRAM alongside the
audio pipeline, FreeRTOS, and networking stack. Three architecture variants
were designed, summarised in Table~\ref{tab:model-evolution}:

\textbf{Full architecture:} 5.0~s audio, 128 Mel bands, filter widths
$[32,64,64]$. The first convolution layer alone produces an intermediate
tensor of approximately 2.05~MB -- more than four times the total SRAM
budget, making this architecture undeployable regardless of quantization.

\textbf{Compact architecture:} Audio duration halved to 2.5~s, Mel bands
halved to 64, filters reduced to $[16,32,32]$, dense layer reduced to 64
units. This reduces the peak intermediate tensor to approximately 257~KB --
still impractical once the log-Mel buffer, FFT buffers, FreeRTOS, I2S DMA,
and Wi-Fi stack overhead are accounted for.

\textbf{Tiny architecture (selected for implementation):} Filters reduced
further to $[8,16,32]$, and \texttt{GlobalAveragePooling2D} used in place of
\texttt{Flatten}+\texttt{Dense} to avoid a large flattened feature vector.
This reduces the peak intermediate tensor to approximately 128~KB, with a
target tensor arena in the 90--212~KB range depending on safety margin.

\begin{table}[H]
    \centering
    \caption{Model architecture design comparison}
    \label{tab:model-evolution}
    \begin{tabular}{lccc}
        \toprule
        \textbf{Parameter} & \textbf{Full} & \textbf{Compact} & \textbf{Tiny} \\
        \midrule
        Audio duration      & 5.0~s          & 2.5~s          & 2.5~s \\
        Mel bands           & 128            & 64             & 64 \\
        Conv filters        & [32, 64, 64]   & [16, 32, 32]   & [8, 16, 32] \\
        Parameters (approx.)& 65{,}000       & 24{,}000       & 9{,}000 \\
        Peak intermediate tensor & $\sim$2.05~MB & $\sim$257~KB & $\sim$128~KB \\
        \bottomrule
    \end{tabular}
\end{table}

\subsection{CNN Architecture}
\label{subsec:cnn-architecture}

The deployed Tiny model is a compact convolutional neural network that
maps a single $64 \times 251 \times 1$ log-Mel spectrogram to a probability
distribution over the five project classes (axe, chainsaw, gunshot,
handsaw, background). It consists of three convolutional blocks followed
by a lightweight classification head, summarised in
Table~\ref{tab:cnn-layers} and illustrated in
Figure~\ref{fig:cnn-block-diagram}.

Each convolutional block applies a $3\times3$ \texttt{Conv2D} layer with
\texttt{same} padding (so the spatial resolution is unchanged by the
convolution itself), followed by \texttt{BatchNormalization} to stabilise
activation statistics during training, a ReLU non-linearity, and a
$2\times2$ \texttt{MaxPooling2D} layer that halves both spatial dimensions.
The three blocks progressively widen the channel dimension
($8 \rightarrow 16 \rightarrow 32$ filters) while shrinking the
frequency-time resolution from $64\times251$ down to $8\times31$, following
the standard CNN design pattern of trading spatial resolution for feature
depth.

Rather than flattening the final $8\times31\times32$ feature map into a
large fully-connected layer -- which would require a
$8\times31\times32=7{,}936$-element weight vector per output neuron --
the network applies \texttt{GlobalAveragePooling2D}, collapsing each of the
32 channels to a single average value. This keeps the head of the network
small (a 32-element vector) regardless of the spectrogram's time-frequency
resolution, which is what makes the \texttt{Tiny} variant deployable within
the ESP32-S3's memory budget (Section~\ref{subsec:model-evolution}). A
\texttt{Dense(64)} layer with ReLU activation and \texttt{Dropout(0.25)}
forms the hidden classification layer, and a final \texttt{Dense(5)} layer
with softmax activation produces the class probabilities. In total the
network has approximately 8,500 trainable parameters (rounded to
$\sim$9,000 in Table~\ref{tab:model-evolution}), the great majority of
which are concentrated in the third convolutional layer and the first
dense layer.

\begin{table}[H]
    \centering
    \caption{Layer-by-layer summary of the deployed Tiny CNN}
    \label{tab:cnn-layers}
    \begin{tabular}{llcc}
        \toprule
        \textbf{Layer} & \textbf{Type} & \textbf{Output Shape} & \textbf{Parameters} \\
        \midrule
        Input        & Log-Mel spectrogram                & $64\times251\times1$ & -- \\
        Conv Block 1 & Conv2D(8, $3\times3$) + BN + ReLU   & $64\times251\times8$ & 112 \\
                     & MaxPool2D($2\times2$)               & $32\times125\times8$ & 0 \\
        Conv Block 2 & Conv2D(16, $3\times3$) + BN + ReLU  & $32\times125\times16$ & 1{,}232 \\
                     & MaxPool2D($2\times2$)               & $16\times62\times16$ & 0 \\
        Conv Block 3 & Conv2D(32, $3\times3$) + BN + ReLU  & $16\times62\times32$ & 4{,}768 \\
                     & MaxPool2D($2\times2$)               & $8\times31\times32$ & 0 \\
        Head         & GlobalAveragePooling2D              & $32$ & 0 \\
                     & Dense(64) + ReLU                    & $64$ & 2{,}112 \\
                     & Dropout($p=0.25$)                   & $64$ & 0 \\
                     & Dense(5) + Softmax                  & $5$ & 325 \\
        \midrule
        \textbf{Total} & & & $\sim$8{,}549 \\
        \bottomrule
    \end{tabular}
\end{table}

\begin{figure}[H]
    \centering
    \includegraphics[width=\linewidth]{cnn_block_diagram.png}
    \caption{Block diagram of the deployed Tiny CNN, showing the three
    convolutional blocks and the classification head}
    \label{fig:cnn-block-diagram}
\end{figure}

\subsection{Quantization Strategy}

Post-training INT8 quantization is applied to each candidate model using
TensorFlow Lite's converter, with a representative calibration dataset used
to fix activation ranges ahead of deployment on target hardware. INT8
quantization is required because the ESP32-S3 has no hardware floating-point
acceleration suited to real-time CNN inference, and materially reduces both
model size and inference latency compared to a float32 model.


\subsection{On-Device Deployment Pipeline}

Figure~\ref{fig:tinyml_pipeline} shows the firmware-level inference pipeline
running on each forest node. The I2S peripheral captures 40{,}000 samples
(2.5~s at 16~kHz) from an INMP441 MEMS microphone into a heap-allocated
\texttt{int16\_t} ring buffer. This buffer is converted into a
$64 \times 251$ log-Mel spectrogram (Hann window, 512-point FFT, 64-band Mel
filterbank, log compression, per-frame normalisation), quantized to INT8
using the input scale and zero-point 
and passed to the TFLite Micro interpreter's \texttt{Invoke()} call. The
raw INT8 output logits are dequantized back to float probabilities
($\text{prob} = (\text{output}_i - \text{zp}) \times \text{scale}$), and the
arg\-max class is taken as the detection result.

\begin{figure}[H]
    \centering
    \makebox[\textwidth][c]{%
        \includegraphics[
            width=1.1\textwidth,
            height=\textheight,
            keepaspectratio
        ]{audiopipeline.pdf}
    }
    \caption{On-device acoustic inference pipeline}
    \label{fig:tinyml_pipeline}
\end{figure}

\subsection{Fire Risk Scoring}
Instead of simple fixed threshold comparisons, each node computes a composite
fire-risk score by multisensory fusion, measuring how much current sensor
readings deviate from normal (calibrated) conditions.
The principle is:
\begin{itemize}
    \item Normal forest conditions = stable temperature, CO$_2$, and
    humidity.
    \item Fire conditions = temperature spikes, CO$_2$ increases
    dramatically, humidity drops.
\end{itemize}
The system calculates a normalized deviation score, and if it exceeds a
threshold, it signals a fire.
\begin{equation}
    \text{FireRisk} = 0.51 \times \Delta CO_2 + 0.37 \times \Delta TEMP +
    0.12 \times \Delta HUMIDITY
\end{equation}
What each part means:
\begin{itemize}
    \item $\Delta CO_2$ = (current CO$_2$ -- Calibrated CO$_2$) / standard
    deviation of CO$_2$
    \item $\Delta T$ = (current temp -- Calibrated temp) / standard deviation
    of temp
    \item $\Delta H$ = (current humidity -- Calibrated humidity) / standard
    deviation of humidity
\end{itemize}
Weights (based on correlation with fire):
\begin{itemize}
    \item 0.51 (51\%) = CO$_2$ (highest correlation: 0.4842)
    \item 0.37 (37\%) = Temperature (correlation: 0.3477)
    \item 0.12 (12\%) = Humidity (lowest correlation: 0.10)
\end{itemize}
Fire likely detected: risk score is above 3.0 threshold.